\begin{examplebox}
\textbf{例1：先读半轴，再求整弦。} 椭圆 $x^2/16+y^2/9=1$ 中 $a=4,b=3$，
故 $\boxed{L=2\cdot9/4=9/2}$。
回验：$c=\sqrt7$，端点 $(\sqrt7,\pm9/4)$ 代入得 $7/16+9/16=1$，两纵坐标之差确为 $9/2$。

\medskip
\textbf{例2：由通径反求参数，并检验存在。}
若椭圆 $b=1$，整通径长为 $\sqrt3$，则 $\sqrt3=2/a$，所以
\[
 \boxed{a=2/\sqrt3=2\sqrt3/3}.
\]
因 $a^2=4/3>1=b^2$，椭圆条件成立；$c^2=1/3$，
端点 $(1/\sqrt3,\pm\sqrt3/2)$ 回代得 $1/4+3/4=1$。

\medskip
\textbf{例3：不同曲线通径相等。}
椭圆离心率 $e=\sqrt2/2$，故 $b^2=a^2(1-e^2)=a^2/2$，通径 $L_E=a$。
若与抛物线 $y^2=2px$ 的通径相等，$a=2p$，所以 $\boxed{p/a=1/2}$。
回验可取 $a=2,b=\sqrt2,p=1$：椭圆端点 $(\sqrt2,\pm1)$、抛物线端点 $(1/2,\pm1)$，两弦长均为 $2$。


\textbf{例4：不用背公式，从焦点坐标重新得出通径。}
求双曲线 $x^2/9-y^2/16=1$ 的通径。
先算 $c=\sqrt{9+16}=5$；通径是 $x=5$，代入
$25/9-y^2/16=1$ 得 $y=\pm16/3$，所以
$\boxed{L=32/3}$。再用 $2b^2/a=2\cdot16/3$ 回验一致。

\medskip
\textbf{例5：题目给的是半通径，不能直接当作整长。}
已知椭圆 $x^2/a^2+y^2/16=1$，其半通径长为 $2$，
则 $16/a=2$，得 $a=8$，整通径实际是 $\boxed{4}$。
检查 $a=8>b=4$，符合长轴沿 $x$ 轴的椭圆条件。

\medskip
\textbf{例6：纵向抛物线，通径应该是水平的。}
抛物线 $x^2=-12y$，因为 $-2p=-12$，有 $p=6$，
焦点 $F=(0,-3)$，通径是水平线 $y=-3$。代入得
$x^2=36$，两端点是 $(-6,-3)$ 和 $(6,-3)$，
因此 $\boxed{L=12=2p}$，不是 $6$。

\medskip
\textbf{例7：看清分母里的 $a$、$b$，不要只按数字大小判断。}
双曲线 $x^2/4-y^2/25=1$，实半轴长 $a=2$，虚半轴长 $b=5$，
故 $\boxed{L=2\cdot25/2=25}$；
不是先看 $25>4$ 就把 $a$ 误认为 $5$。

\end{examplebox}
\clearpage
